\documentclass[11pt]{article}
\usepackage[a4paper,margin=27mm]{geometry}
\usepackage[T1]{fontenc}
\usepackage{lmodern,microtype,amsmath,amssymb,amsthm,mathtools,enumitem}
\usepackage{xcolor}
\usepackage[colorlinks=true,linkcolor=blue,citecolor=blue,urlcolor=blue]{hyperref}
\hypersetup{pdftitle={Coherence and Dominance under General Rankings},hypertexnames=false}
\newtheorem{theorem}{Theorem}[section]
\newtheorem{lemma}[theorem]{Lemma}
\newtheorem{proposition}[theorem]{Proposition}
\newtheorem{corollary}[theorem]{Corollary}
\theoremstyle{definition}
\newtheorem{definition}[theorem]{Definition}
\newtheorem{example}[theorem]{Example}
\theoremstyle{remark}
\newtheorem{remark}[theorem]{Remark}
\newcommand{\M}{\mathcal M}
\newcommand{\C}{\mathcal C}
\newcommand{\I}{\mathcal I}
\newcommand{\U}{\mathcal U}
\newcommand{\Lf}{\mathfrak L}
\newcommand{\supp}{\operatorname{supp}}
\setlist{itemsep=1pt,topsep=3pt}

\title{Dominance Generalised}
\author{Jason Konek}
\date{October 2026}

\begin{document}
\maketitle

\section{Introduction}


Dominance arguments are commonly used to justify the claim that some axioms or other are requirements of rationality. For example, De~Finetti argued that an agent's forecasts or estimates of bounded random variables ought to be probabilistic expectations \cite[\S1.6, \S3.3]{deFinetti1974}. This is necessary and sufficient for any finite collection of those estimates to avoid uniform dominance when evaluated using a loss function known as the Brier score. A finite set of estimates is uniformly dominated if there is some alternative set of forecasts/estimates that is guaranteed to get a smaller penalty (according to the loss function). Dominance seems to indicate a structural defect along whatever dimension of evaluation the loss function tracks. If the Brier score tracks a dimension of evaluation that is of cardinal importance to rationality, then forecasts/estimates that are dominated are plausibly rationally impermissible. Hence non-probabilistic estimates are rationally impermissible. Probabilistic estimates, in contrast, are rationally permissible (or at least they pass every test the dominance argument considers).


A number of authors have provided dominance arguments for coherence norms like probabilism~\cite{Joyce2009,Predd2009,SSK2009,Pettigrew2022,Nielsen2023}, norms governing updating like conditionalization~\cite{Briggs2020,Nielsen2021}, and deference norms like the Principal Principle~\cite[Chapter~10]{Pettigrew2016}. As elegant and insightful as these arguments are, they also muddle together some things that we really ought to keep distinct.




\section{Setting}\label{sec:setting}

\begin{definition}[Rankings]\label{def:rankings}
Let $\Omega$ be a nonempty set of \emph{worlds} and $\M$ a nonempty set of \emph{models}. A \emph{ranking} $L$ consists of a nonempty set $\M_L\subseteq\M$, whose members are the models \emph{evaluable} by $L$, and, for each world $\omega$, a binary relation $\preceq_\omega$ on $\M_L$. We read $m\preceq_\omega m'$ as: $m$ performs at least as well as $m'$ at $\omega$. Models $m$ and $m'$ are \emph{tied at $\omega$} if $m\preceq_\omega m'$ and $m'\preceq_\omega m$. No additional property of these binary relations is assumed. For a fixed ranking $L$, ``model'' means a member of $\M_L$.
\end{definition}

\begin{definition}[Dominance]\label{def:dominance}
Fix a ranking $L$ and models $c,m\in\M_L$.
\begin{itemize}[leftmargin=*]
\item $c$ \emph{weakly dominates} $m$ if $c\preceq_\omega m$ at every world.
\item $m$ and $c$ are \emph{tied at every world} if each weakly dominates the other.
\item $c$ \emph{dominates} $m$ if $(\forall\omega\in\Omega)\ c\preceq_\omega m$ and $(\exists\omega\in\Omega)\ m\not\preceq_\omega c$.
\end{itemize}
A model is \emph{undominated} if no model dominates it. $\U_L$ is the set of undominated models.
\end{definition}

\begin{definition}[Probability mass functions]\label{def:priors}
For a finite nonempty $F\subseteq\Omega$, put $\Delta_F\coloneqq\{p\in[0,1]^F:\sum_{\omega\in F}p_\omega=1\}$, with the Euclidean topology. Its elements are probability mass functions (pmfs) on $F$; for $A\subseteq\Omega$, put $p(A)=\sum_{\omega\in A\cap F}p_\omega$. The support of a pmf $p$ is $\supp p\coloneqq\{\omega\in F:p_\omega>0\}$. For $\omega\in F$, let $\delta_\omega\in\Delta_F$ have coordinate $1$ at $\omega$ and $0$ elsewhere. When $\Omega$ is finite, $\Delta\coloneqq\Delta_\Omega$.
\end{definition}

\begin{definition}[Coherent models]\label{def:coherence-property}
A class $\C\subseteq\M$ of \emph{coherent} models is fixed in advance, independently of the rankings. Models not in $\C$ are \emph{incoherent}. For a ranking $L$, the \emph{undominated coherent models} of $L$ are
\[
\C_L\coloneqq\{c\in\C\cap\M_L:\ \text{no } m\in\M_L \text{ dominates } c\}=\C\cap\U_L.
\]
\end{definition}

\begin{definition}[Closure under ties]\label{def:tie-closure}
Fix a ranking $L$. A set $\mathcal A\subseteq\M_L$ is \emph{closed under ties} if
\[
\forall a\in\mathcal A\ \forall m\in\M_L:\quad
\bigl[(\forall\omega\in\Omega)\ (a\preceq_\omega m\ \textcolor{red}{\text{and}}\ m\preceq_\omega a)\bigr]\ \Longrightarrow\ m\in\mathcal A.
\] A ranking $L$ is \emph{closed under ties} if $\C_L$ is.
\end{definition}

\begin{definition}[Evaluability and coherent admissibility]\label{def:family-assumptions}
Let $\Lf$ be a family of rankings.
\begin{itemize}[leftmargin=*]
\item \emph{Evaluability}: $\M=\bigcup_{L\in\Lf}\M_L$.
\item \emph{Coherent admissibility}: $\C=\bigcup_{L\in\Lf}\C_L$.
\end{itemize}
\end{definition}

\begin{lemma}\label{lem:ties}
Let a model $a$ weakly dominate a model $m$. Then $a$ dominates $m$ if and only if $m$ does not weakly dominate $a$.
\begin{enumerate}[label=(\alph*)]
\item If $m$ is undominated, then $m$ and $a$ are tied at every world.
\item If $L$ is closed under ties, $a\in\C_L$ and $m\notin\C_L$, then $a$ dominates $m$.
\item More generally, if $\mathcal A$ is closed under ties, $a\in\mathcal A$ and $m\notin\mathcal A$, then $a$ dominates $m$.
\end{enumerate}
\end{lemma}
\begin{proof}
The first claim is the definition of dominance. (a) Otherwise $a$ would dominate $m$. (c) Since $m\notin\mathcal A$, $m$ is not tied at every world with $a$. (b) This is (c) with $\mathcal A\coloneqq\C_L$.
\end{proof}

\section{Selection by a countable family of assignments}\label{sec:general}

\begin{lemma}[Knaster--Kuratowski--Mazurkiewicz {\cite{KKM1929}}]\label{lem:kkm}
Let $F$ be finite and nonempty, and let $(A_\omega)_{\omega\in F}$ be closed subsets of $\Delta_F$ such that every $p\in\Delta_F$ lies in $A_\omega$ for some $\omega\in\supp p$. Then $\bigcap_{\omega\in F}A_\omega\ne\varnothing$.
\end{lemma}
\begin{proof}
Each $A_\omega$ is nonempty, since by hypothesis we always have
\[
\delta_\omega\in\bigcup_{\omega^\prime\in\supp \delta_\omega}A_{\omega^\prime}=A_\omega.
\]
Let $d_\omega(p)\coloneqq\min_{q\in A_\omega}\lVert p-q \rVert$ and $s(p)\coloneqq\sum_{\omega\in F} d_\omega(p)$. The map $\Phi(p)_\omega\coloneqq(p_\omega+d_\omega(p))/(1+s(p))$ is continuous. Moreover, $\Phi$ maps each $p\in\Delta_F$ to some point in $\Delta_F$ (since each coordinate $\Phi(p)_\omega$ of $\Phi(p)$ is nonnegative and these coordinates sum to 1). Brouwer's fixed-point theorem then guarantees that there is some $p\in\Delta_F$ with $\Phi(p)=p$. This holds just in case $p_\omega\cdot s(p)=d_\omega(p)$ for all $\omega\in F$. By assumption, there is some $\omega\in\supp p$ with $p\in A_\omega$. Choose some such $\omega$. Then $d_\omega(p)=0$, so $p_\omega\cdot s(p)=0$. Since $p_\omega>0$, $s(p)=0$ and hence $p\in\bigcap_{\omega\in F}A_\omega$.
\end{proof}

\subsection{Finitely many worlds}

\begingroup\color{red}
\begin{definition}[Countable selection, finitely many worlds]\label{def:selection-finite}
Let $\Omega$ be finite, $L$ a ranking and $\I\subseteq\M_L$. Then $L$ satisfies \emph{countable selection on $\I$} if there are models $c^j_p\in\M_L$, for $j\in\mathbb N=\{1,2,\ldots\}$ and $p\in\Delta$, such that for every $x\in\I$ there exist an index $j\in\mathbb N$ and closed sets $A_\omega(x)\subseteq\Delta$, $\omega\in\Omega$, satisfying
\begin{align}
&\forall\omega\in\Omega\ \forall p\in A_\omega(x):\ c^j_p\preceq_\omega x;\tag{F1}\label{eq:f-comparison}\\
&\forall p\in\Delta\ \exists\omega\in\Omega:\ p_\omega>0\ \text{and}\ p\in A_\omega(x).\tag{F2}\label{eq:f-support}
\end{align}
The models $c^j_p$ are chosen before $x$; only $j$ and the sets $A_\omega(x)$ may depend on $x$. \emph{Single-assignment selection} is the same condition with $\mathbb N$ replaced by $\{1\}$.
\end{definition}

\paragraph{Informal reading.} Think of each index $j$ as an \emph{adviser}, fixed before any rival is considered, who recommends the model $c^j_p$ to anyone whose probabilities over the worlds are $p$. Given a rival $x$, one picks an adviser. $A_\omega(x)$ is a closed set of probabilities whose recommendations do at least as well as $x$ at the world $\omega$ (F1); they need not do as well as $x$ at any other world. (F2) says that, for every $p$, the recommendation for $p$ does at least as well as $x$ at some world to which $p$ gives positive probability. Closedness means that this cannot fail only in the limit. Theorem~\ref{thm:finite} then finds a probability whose recommendation does at least as well as $x$ at every world, that is, a weak dominator of $x$. Whether it actually dominates $x$ is settled by the condition on ties in Section~\ref{sec:conclusion}. The advisers are countably many and fixed in advance; only the choice of adviser depends on $x$.

\begin{theorem}\label{thm:finite}
Let $\Omega$ be finite and $L$ a ranking. Let models $c^j_p$ be as in Definition~\ref{def:selection-finite}, let $x\in\M_L$, and let an index $j$ and closed sets $A_\omega(x)\subseteq\Delta$ satisfy (F1) and (F2). Then $\bigcap_{\omega\in\Omega}A_\omega(x)\ne\varnothing$, and $c^j_p$ weakly dominates $x$ for every $p$ in this intersection. If no such $c^j_p$ is tied with $x$ at every world, every such $c^j_p$ dominates $x$.
\end{theorem}
\begin{proof}
By (F2), the closed sets $A_\omega(x)$ satisfy the hypothesis of Lemma~\ref{lem:kkm} with $F\coloneqq\Omega$, so they have a common point. For every $p$ in their intersection, (F1) gives $c^j_p\preceq_\omega x$ at every world, so $c^j_p$ weakly dominates $x$. If, moreover, $c^j_p$ and $x$ are not tied at every world, $c^j_p$ dominates $x$.
\end{proof}
\endgroup

\subsection{Arbitrary sets of worlds}

{\color{red}With infinitely many worlds there may be no probability on $\Omega$ to index the assignments by. Instead, the assigned models are indexed by a compact space, and finite sets of worlds are handled one at a time by maps from the probability simplices on those sets into that space.}

\begin{definition}[Countable selection]\label{def:selection}\label{def:family}
{\color{red}Let $\Omega$ be any nonempty set,} $L$ a ranking and $\I\subseteq\M_L$. Write $\mathcal F(\Omega)$ for the set of all finite subsets of $\Omega$, including $\varnothing$. A topological space is \emph{compact} if every cover by open sets has a finite subcover. A subset is \emph{closed} if its complement is open. A map between topological spaces is \emph{continuous} if the inverse image of each open set is open.

$L$ satisfies \emph{countable selection on $\I$} if there is a sequence
\[
\bigl(\Theta_j,b^j,F_j,(t_{j,G})_{G\in\mathcal F(\Omega)}\bigr)_{j\in\mathbb N},\qquad \mathbb N=\{1,2,\ldots\},
\]
with the following properties for every $j\in\mathbb N$:
\begin{enumerate}[label=(\roman*)]
\item a nonempty compact topological space $\Theta_j$ and a function $b^j:\Theta_j\to\M_L$, with $b^j_\theta\coloneqq b^j(\theta)$;
\item a function $F_j:\mathcal F(\Omega)\to\mathcal F(\Omega)\setminus\{\varnothing\}$ with $G\subseteq F_j(G)$ for every $G\in\mathcal F(\Omega)$;
\item for every $G\in\mathcal F(\Omega)$, a function $t_{j,G}:\Delta_{F_j(G)}\to\Theta_j$;
\end{enumerate}
such that for every $x\in\I$ there exist an index $j\in\mathbb N$, closed sets $K_\omega(x)\subseteq\Theta_j$ for all $\omega\in\Omega$, and closed sets $A^G_\omega(x)\subseteq\Delta_{F_j(G)}$ for all $G\in\mathcal F(\Omega)$ and $\omega\in F_j(G)$, satisfying
\begin{align}
&\forall\omega\in\Omega\ \forall\theta\in K_\omega(x):
 b^j_\theta\preceq_\omega x;\tag{W1}\label{eq:comparison}\\
&\forall G\in\mathcal F(\Omega)\ \forall\omega\in F_j(G):
 t_{j,G}(A^G_\omega(x))\subseteq K_\omega(x);\tag{W2}\label{eq:compatibility}\\
&\forall G\in\mathcal F(\Omega)\ \forall p\in\Delta_{F_j(G)}\
 \exists\omega\in F_j(G):
 p_\omega>0\ \text{and}\ p\in A^G_\omega(x).\tag{W3}\label{eq:support}
\end{align}
Condition (W1) requires the stated comparisons; (W2) preserves them under the fixed maps; (W3) requires one at a world given positive probability.

The spaces and functions in (i)--(iii) are chosen before $x$; only $j$, $K_\omega(x)$ and $A^G_\omega(x)$ may depend on $x$. No continuity of $b^j$ or $t_{j,G}$ is required. \emph{Single-assignment selection} is the same condition with the index set $\mathbb N$ replaced by $\{1\}$.
\end{definition}

{\color{red}\paragraph{Informal reading.} The adviser $j$ now has a compact stock $\Theta_j$ of recommendations, recommendation $\theta$ being the model $b^j_\theta$ (i). Asked to weigh a finite set $G$ of worlds, the adviser enlarges it to a finite set $F_j(G)\supseteq G$ (ii) and, for each probability $p$ on $F_j(G)$, makes the recommendation $t_{j,G}(p)$ (iii). For a rival $x$, $K_\omega(x)$ is a closed set of recommendations that do at least as well as $x$ at $\omega$ (W1). $A^G_\omega(x)$ is a closed set of probabilities on $F_j(G)$ whose recommendations lie in $K_\omega(x)$ (W2). (W3) is (F2) for each finite set of worlds. Theorem~\ref{thm:general} applies Theorem~\ref{thm:finite}'s argument to each finite set and uses compactness of $\Theta_j$ to combine the answers into one recommendation that does at least as well as $x$ at every world. Proposition~\ref{prop:finite-reduction} shows that for finite $\Omega$ nothing is lost by using Definition~\ref{def:selection-finite}.}

\begingroup\color{red}
\begin{proposition}\label{prop:finite-reduction}
Let $\Omega$ be finite. Then $L$ satisfies countable selection on $\I$ in the sense of Definition~\ref{def:selection} if and only if it does in the sense of Definition~\ref{def:selection-finite}, and then the assigned models can be taken to be the same.
\end{proposition}
\begin{proof}
Suppose Definition~\ref{def:selection-finite} holds. For every $j$ and finite $G$ put $\Theta_j\coloneqq\Delta$, which is compact, $b^j_p\coloneqq c^j_p$, $F_j(G)\coloneqq\Omega$ and $t_{j,G}$ the identity on $\Delta$. For each $x\in\I$ take its index $j$ and put $K_\omega(x)\coloneqq A_\omega(x)$ and $A^G_\omega(x)\coloneqq A_\omega(x)$. Then (F1) gives (W1), (W2) is trivial, and (F2) gives (W3).

Conversely, suppose Definition~\ref{def:selection} holds. Since $\Omega\subseteq F_j(\Omega)\subseteq\Omega$, $F_j(\Omega)=\Omega$. Put $c^j_p\coloneqq b^j_{t_{j,\Omega}(p)}$ for $p\in\Delta$, and, for $x\in\I$ with its index $j$, $A_\omega(x)\coloneqq A^\Omega_\omega(x)$. These sets are closed. (W2) and (W1) give (F1), and (W3) with $G\coloneqq\Omega$ gives (F2). Every $c^j_p$ is an assigned model $b^j_\theta$.
\end{proof}
\endgroup


\begin{theorem}[Countable selection]\label{thm:general}
Let $\Omega$ be any nonempty set and $L$ any ranking. Fix spaces and functions as in Definition~\ref{def:selection}(i)--(iii), a model $x\in\M_L$, an index $j$, and closed sets $K_\omega(x)$ and $A^G_\omega(x)$ satisfying (W1)--(W3). Then $\bigcap_{\omega\in\Omega}K_\omega(x)\ne\varnothing$, and for every $\theta$ in this intersection $b^j_\theta$ weakly dominates $x$. {\color{red}If no such $b^j_\theta$ is tied with $x$ at every world, every such $b^j_\theta$ dominates $x$.}
\end{theorem}
\begin{proof}
Let $G\subseteq\Omega$ be finite and put $F\coloneqq F_j(G)$. By (W3), the closed sets $A^G_\omega(x)$, $\omega\in F$, satisfy the hypothesis of Lemma~\ref{lem:kkm}. So some $p_G$ lies in all of them. By (W2), $t_{j,G}(p_G)\in K_\omega(x)$ for every $\omega\in F$, and in particular for every $\omega\in G$, since $G\subseteq F$. Taking $G\coloneqq\{\omega_1,\dots,\omega_k\}$ for arbitrary worlds $\omega_1,\dots,\omega_k$ shows that every finite subfamily of the closed sets $K_\omega(x)$, $\omega\in\Omega$, has nonempty intersection. It follows that $\bigcap_{\omega\in\Omega}K_\omega(x)\ne\varnothing$. Otherwise the open sets $\Theta_j\setminus K_\omega(x)$, $\omega\in\Omega$, would cover $\Theta_j$. By compactness, finitely many of them, $\Theta_j\setminus K_{\omega_1}(x),\dots,\Theta_j\setminus K_{\omega_k}(x)$, would cover $\Theta_j$. But then $K_{\omega_1}(x)\cap\dots\cap K_{\omega_k}(x)=\varnothing$, a contradiction.

Let $\theta\in\bigcap_{\omega\in\Omega}K_\omega(x)$. By (W1), $b^j_\theta\preceq_\omega x$ at every world, so $b^j_\theta$ weakly dominates $x$. {\color{red}If, moreover, $b^j_\theta$ and $x$ are not tied at every world, $b^j_\theta$ dominates $x$.}
\end{proof}

\begin{corollary}\label{cor:inverse}
Fix the spaces and functions in Definition~\ref{def:selection}(i)--(iii), and suppose every $t_{j,G}$ is continuous. Let $x$ be a model, $j$ an index, and $K_\omega(x)\subseteq\Theta_j$, $\omega\in\Omega$, closed sets satisfying (W1), such that for every finite $G\subseteq\Omega$ and every $p\in\Delta_{F_j(G)}$ there is some $\omega\in\supp p$ with $t_{j,G}(p)\in K_\omega(x)$. Then $K_\omega(x)$ and $A^G_\omega(x)\coloneqq t_{j,G}^{-1}(K_\omega(x))$ satisfy (W1)--(W3) for $x$ at $j$.
\end{corollary}
\begin{proof}
The sets $A^G_\omega(x)$ are closed by continuity, (W2) holds by construction, and (W3) is the displayed hypothesis.
\end{proof}


\begingroup\color{red}
\begin{lemma}\label{lem:no-tie}
Fix assigned models as in Definition~\ref{def:selection-finite} or Definition~\ref{def:selection}, and let $x\in\M_L$ be a model with which no assigned model is tied at every world. Then, for every index $j$ and all closed sets satisfying (F1)--(F2), respectively (W1)--(W3),
\begin{align}
&\forall p\in\bigcap_{\omega\in\Omega}A_\omega(x)\ \exists\omega\in\Omega:\quad x\not\preceq_\omega c^j_p,\tag{F3}\label{eq:f-notie}\\
&\forall\theta\in\bigcap_{\omega\in\Omega}K_\omega(x)\ \exists\omega\in\Omega:\quad x\not\preceq_\omega b^j_\theta,\tag{W4}\label{eq:notie}
\end{align}
respectively, and every assigned model at a common point of these sets dominates $x$.
\end{lemma}
\begin{proof}
At a common point, (F1), respectively (W1), makes the assigned model weakly dominate $x$. As the two are not tied at every world, $x$ does not weakly dominate it, which is (F3), respectively (W4); so it dominates $x$.
\end{proof}

\begin{lemma}\label{lem:tie-sufficient}
If every assigned model lies in $\C_L$ and $\C_L$ is closed under ties, then no assigned model is tied at every world with a model outside $\C_L$, in particular with an incoherent model.
\end{lemma}
\begin{proof}
If an assigned model $b\in\C_L$ were tied at every world with $x$, closure under ties would give $x\in\C_L$. Incoherent models are not in $\C_L\subseteq\C$.
\end{proof}

\begin{corollary}\label{cor:coherent-assigned}
Let $L$ satisfy countable selection on $\I\subseteq\M_L$ with every assigned model in $\C_L$.
\begin{enumerate}[label=(\alph*)]
\item If no assigned model is tied at every world with a model in $\I\setminus\C_L$, every $x\in\I\setminus\C_L$ is dominated by a member of $\C_L$.
\item If $\C_L$ is closed under ties, every $x\in\I\setminus\C_L$ is dominated by a member of $\C_L$.
\end{enumerate}
\end{corollary}
\begin{proof}
(a) Theorem~\ref{thm:general} and Lemma~\ref{lem:no-tie} give a dominating assigned model, which lies in $\C_L$. (b) By Lemma~\ref{lem:tie-sufficient}, (a) applies.
\end{proof}
\endgroup

\section{The dominance argument}\label{sec:conclusion}

\begin{definition}[Dominance principle]\label{def:permissibility}
Let $\Lf$ be a family of rankings. For $m\in\M$, define:
\begin{align*}
m\text{ is irrational}&\quad\Longleftrightarrow\quad (\forall L\in\Lf)\bigl[m\in\M_L\Rightarrow m\notin\U_L\bigr];\\
m\text{ is rationally permissible}&\quad\Longleftrightarrow\quad (\exists L\in\Lf)\ m\in\U_L.
\end{align*}
\end{definition}

\begin{lemma}\label{lem:complement}
A model is irrational if and only if it is not rationally permissible. Under evaluability, an irrational model is dominated under at least one ranking in $\Lf$.
\end{lemma}
\begin{proof}
A model fails to be irrational exactly when some ranking in $\Lf$ that evaluates it leaves it undominated, which is rational permissibility. For the second sentence, by evaluability some $L\in\Lf$ evaluates an irrational model, and irrationality gives a dominator under that $L$.
\end{proof}

\begin{definition}[Assumptions on the rankings]\label{def:argument-assumptions}
{\color{red}Conditions (S) and (T)} are conditions on a single ranking $L$; the family $\Lf$ satisfies them if every $L\in\Lf$ does. The dominance argument uses the following conditions.
\begin{itemize}[leftmargin=3.4em,labelsep=0.5em]
\item[\textbf{(E)}] \emph{Evaluability}: $\M=\bigcup_{L\in\Lf}\M_L$. No model escapes assessment.
\item[\textbf{(A)}] \emph{Coherent admissibility}: $\C=\bigcup_{L\in\Lf}\C_L$. Each coherent model is undominated under at least one ranking that evaluates it.
\item[\textbf{{\color{red}(S)}}] \emph{Countable selection}: $L$ satisfies countable selection on $\M_L\setminus\C$ (Definition~\ref{def:selection}{\color{red}; for finite $\Omega$, equivalently Definition~\ref{def:selection-finite}, by Proposition~\ref{prop:finite-reduction}}). The functions are fixed before the incoherent model is considered.
{\color{red}\item[\textbf{(T)}] \emph{Exclusion of ties}: the assigned models in (S) can be chosen so that, in addition, none of them is tied at every world with an incoherent model in $\M_L$. No incoherent model performs exactly as well as an assigned model at every world.}
\end{itemize}
\end{definition}

\begin{theorem}[The dominance argument]\label{thm:argument}
Let $\Lf$ be a family of rankings.
\begin{enumerate}[label=(\alph*)]
\item If a ranking $L$ satisfies {\color{red}(S)} and {\color{red}(T)}, then every incoherent model in $\M_L$ is dominated under $L$ by an assigned model, and $\U_L=\C_L$.
\item If $\Lf$ satisfies {\color{red}(S)} and {\color{red}(T)}, every incoherent model is irrational.
\item If $\Lf$ satisfies (A), every coherent model is rationally permissible.
\item If $\Lf$ satisfies (A), {\color{red}(S)} and {\color{red}(T)}, the rationally permissible models are exactly the coherent ones, and the irrational models exactly the incoherent ones.
\item If $\Lf$ satisfies (E), {\color{red}(S)} and {\color{red}(T)}, every incoherent model is dominated under at least one ranking in $\Lf$.
\end{enumerate}
\end{theorem}
\begin{proof}
\leavevmode
\begin{enumerate}[label=\arabic*.,leftmargin=*]
\item Let $L$ satisfy {\color{red}(S)} and {\color{red}(T)}, and let $x\in\M_L$ be incoherent. {\color{red}By (S) and (T), $x$ has an index $j$ and closed sets satisfying (W1)--(W3), and no assigned model is tied with $x$ at every world. By Theorem~\ref{thm:general} and Lemma~\ref{lem:no-tie}, $x$ is dominated by an assigned model $b^j_\theta$.} Now let $m\in\U_L$. Then $m$ is not dominated, so it is not incoherent, and $m\in\C\cap\U_L=\C_L$. Conversely $\C_L\subseteq\U_L$ by definition. This is (a).
\item By (a), every incoherent model is dominated under every ranking in $\Lf$ that evaluates it; that is, it is irrational. This is (b).
\item By (A), each coherent model $c$ belongs to $\C_L$ for some $L\in\Lf$. Then $L$ evaluates $c$ and leaves it undominated, so $c$ is rationally permissible. This is (c).
\item By (b), (c) and Lemma~\ref{lem:complement}, no incoherent model is rationally permissible and no coherent model is irrational. Every model is coherent or incoherent, so (d) follows.
\item By (E), each incoherent model is evaluated by some $L\in\Lf$, and by (a) it is dominated under $L$. This is (e).\qedhere
\end{enumerate}
\end{proof}

\begin{corollary}[Coherent dominators]\label{cor:coherent-dominators}
Let a ranking $L$ satisfy countable selection on $\M_L\setminus\C$ with $b^j_\theta\in\C_L$ for every $j\in\mathbb N$ and $\theta\in\Theta_j$. If {\color{red}no assigned model is tied at every world with an incoherent model}, or if $L$ is closed under ties, then every incoherent model in $\M_L$ is dominated by a member of $\C_L$, and $\U_L=\C_L$. Consequently, if every $L\in\Lf$ is of this kind, parts (b)--(e) of Theorem~\ref{thm:argument} hold under their stated hypotheses on (A) and (E), and the dominators in (e) may be taken in $\C_L$.
\end{corollary}
\begin{proof}
If $L$ is closed under ties, Lemma~\ref{lem:tie-sufficient} gives {\color{red}(T)}. Corollary~\ref{cor:coherent-assigned}(a), with $\I\coloneqq\M_L\setminus\C$, gives the dominators in $\C_L$, and Theorem~\ref{thm:argument}(a) gives $\U_L=\C_L$. The rest is Theorem~\ref{thm:argument}.
\end{proof}

\section{Countable selection and a single assignment}\label{sec:scope}

\begin{proposition}\label{prop:comparison}
Let $\I\subseteq\M_L$.
\begin{enumerate}[label=(\alph*)]
\item Single-assignment selection on $\I$ implies countable selection on $\I${\color{red}, with the same assigned models}.
\item If there is a sequence $(d_j)_{j\ge1}$ in $\M_L$ such that for every $x\in\I$ some $d_j$ dominates $x$, {\color{red}then $L$ satisfies countable selection on $\I$ with the $d_j$ as assigned models.}
\end{enumerate}
\end{proposition}
\begin{proof}
(a) Repeat the assignment. (b) For each $j$ take a one-point space $\Theta_j$, assign $d_j$, put $F_j(G)=G\cup\{\omega_0\}$ for a fixed $\omega_0\in\Omega$, and let every $t_{j,G}$ be constant. If $d_j$ dominates $x$, take $K_\omega(x)=\Theta_j$ and $A^G_\omega(x)=\Delta_{F_j(G)}$ for every $G$ and $\omega$. {\color{red}Conditions (W1)--(W3) hold.}
\end{proof}

\begin{example}\label{ex:two-pairs}
Let $\Omega\coloneqq\{\omega\}$ and $\M\coloneqq\{c_1,x_1,c_2,x_2\}$, four distinct models, with $\C\coloneqq\{c_1,c_2\}$. Let $L$ evaluate every model, with $\preceq_\omega$ consisting of $m\preceq_\omega m$ for every $m$, together with $c_1\preceq_\omega x_1$ and $c_2\preceq_\omega x_2$. Each $c_i$ is undominated, ties occur only between a model and itself, and so $\C_L=\{c_1,c_2\}$ is closed under ties. {\color{red}Countable selection on $\{x_1,x_2\}$ holds, with no assigned model tied at every world with $x_1$ or $x_2$, by} repeating the following two choices for odd and even $j$, respectively: $\Theta_i$ a single point $\theta$, $b^i_\theta\coloneqq c_i$, $F_i(G)\coloneqq\Omega$, $t_{i,G}(p)\coloneqq\theta$, and, for $x_i$, index $i$ with $K_\omega(x_i)\coloneqq\Theta_i$ and $A^G_\omega(x_i)\coloneqq\Delta_\Omega$. Single-assignment selection on $\{x_1,x_2\}$ fails. For $G\coloneqq\Omega$, the map $t_{1,G}$ sends the only pmf $\delta_\omega$ to one parameter $\theta_0$. By (W3) and (W2), $\theta_0\in K_\omega(x_1)\cap K_\omega(x_2)$, so by (W1) $b^1_{\theta_0}\preceq_\omega x_1$ and $b^1_{\theta_0}\preceq_\omega x_2$. But no model is at least as good as both $x_1$ and $x_2$. So countable selection is strictly weaker than single-assignment selection, even on the incoherent models and with every assigned model in $\C_L$.
\end{example}

\begin{example}\label{ex:uncountable}
Let $\Omega\coloneqq\{\omega\}$, let $T$ be uncountable, and let $\M\coloneqq\{c_t,x_t:t\in T\}$, all distinct, with $\C\coloneqq\{c_t:t\in T\}$. Let $L$ evaluate every model, with $\preceq_\omega$ consisting of $m\preceq_\omega m$ for every $m$, together with $c_t\preceq_\omega x_t$ for every $t$. Each $x_t$ is dominated by $c_t$. Countable selection on $\{x_t:t\in T\}$ fails. Given any spaces and functions as in Definition~\ref{def:selection}(i)--(iii), let $\theta_j\coloneqq t_{j,\Omega}(\delta_\omega)$. As in Example~\ref{ex:two-pairs}, if (W1)--(W3) hold for $x_t$ at index $j$, then $b^j_{\theta_j}\preceq_\omega x_t$, so $b^j_{\theta_j}\in\{c_t,x_t\}$. This determines $t$, so each index serves at most one $x_t$, and countably many indices serve only countably many. So dominance of every incoherent model does not imply countable selection.
\end{example}



\section{Applications to accuracy}\label{sec:short-applications}
For a loss $L:\M_L\times\Omega\to[0,\infty]$, put $m\preceq_\omega m'$ exactly when $L(m,\omega)\le L(m',\omega)$. For finite $\Omega$, put $E_pL(m,\cdot)=\sum_{\omega\in\Omega}p_\omega L(m,\omega)$, with $0\cdot\infty=0$. A \emph{Bayes model for $p$} is a model $c$ such that $E_pL(c,\cdot)\le E_pL(m,\cdot)$ for every $m\in\M_L$. A function $f$ into $[0,\infty]$ is \emph{lower semicontinuous} if $\{z:f(z)\le a\}$ is closed for every real $a$. The conclusions below concern dominance as defined here: improvement at some world, without worsening at any. Improvement at every world is a stronger conclusion. {\color{red}Each proposition below verifies {\color{red}(S)} and {\color{red}(T)} by the following lemma, which uses a single assignment.}

\begingroup\color{red}
\begin{lemma}\label{lem:loss-selection}
Let $L$ be given by a loss on $\M_L$. Let $\Theta$ be a compact space and $\theta\mapsto c_\theta$ a map from $\Theta$ into $\C_L$ such that $\theta\mapsto L(c_\theta,\omega)$ is lower semicontinuous, as a map into $[0,\infty]$, at every world. Suppose that for every finite nonempty $F\subseteq\Omega$ there is a continuous map $\theta_F:\Delta_F\to\Theta$ such that, for every $m\in\M_L$ and every $p\in\Delta_F$, some $\omega\in\supp p$ has $L(c_{\theta_F(p)},\omega)\le L(m,\omega)$. If $\C_L$ is closed under ties, {\color{red}then $L$ satisfies (S) and (T) with a single assignment}, every model not in $\C_L$ is dominated by some $c_\theta$, and $\U_L=\C_L$.
\end{lemma}
\begin{proof}
Fix a world $\omega_0$. For every $j$ put $\Theta_j\coloneqq\Theta$ and $b^j_\theta\coloneqq c_\theta$; for every finite $G\subseteq\Omega$ put $F_j(G)\coloneqq G$ if $G\ne\varnothing$, $F_j(\varnothing)\coloneqq\{\omega_0\}$, and $t_{j,G}\coloneqq\theta_{F_j(G)}$. This is a single assignment with every assigned model in $\C_L\subseteq\M_L$. For a model $m$ put $K_\omega(m)\coloneqq\{\theta:L(c_\theta,\omega)\le L(m,\omega)\}$, which is closed by lower semicontinuity (and equals $\Theta$ when $L(m,\omega)=\infty$). (W1) holds by definition, and the displayed hypothesis is the support condition of Corollary~\ref{cor:inverse}. Explicitly, put $A^G_\omega(m)=t_{j,G}^{-1}(K_\omega(m))$. These closed sets satisfy (W1)--(W3) for every $m\in\M_L$, by Corollary~\ref{cor:inverse}; {\color{red}by Lemma~\ref{lem:tie-sufficient}, no assigned model is tied at every world with a model outside $\C_L$. This gives (S) and (T);} Theorem~\ref{thm:general} {\color{red}and Lemma~\ref{lem:no-tie} give} a dominating $c_\theta$ for every model not in $\C_L$; and Theorem~\ref{thm:argument}(a) gives $\U_L=\C_L$.
\end{proof}
\endgroup

\subsection{Proper scores}
\begin{proposition}\label{prop:short-proper}
Let $\Omega$ be finite and let $p\mapsto c_p$ map $\Delta$ onto $\C\subseteq\M_L$. Suppose each $c_p$ is the unique Bayes model for $p$, and $p\mapsto L(c_p,\omega)$ is lower semicontinuous at every world. Then the undominated models are exactly the coherent models, and every other model is dominated by a coherent one.
\end{proposition}
\begin{proof}
If $\M_L$ is a singleton, the conclusion is immediate. Otherwise uniqueness makes each Bayes expected loss finite: if the minimum were infinite, every model would attain it. {\color{red}A weak dominator of $c_p$ is also Bayes for $p$, hence equals $c_p$; in particular, so is any model tied with $c_p$ at every world. Thus $\C_L=\C$ and $\C_L$ is closed under ties. Bayes optimality gives, for each $m$, a world in $\supp p$ where $L(c_p,\omega)\le L(m,\omega)$: strict improvement throughout the support would strictly reduce the finite expected loss. Apply Lemma~\ref{lem:loss-selection} with $\Theta\coloneqq\Delta$ and, for each nonempty $F\subseteq\Omega$, $\theta_F(p)_\omega=p_\omega$ for $\omega\in F$ and $\theta_F(p)_\omega=0$ otherwise; the support of $\theta_F(p)$ is $\supp p$.}
\end{proof}
Taking $c_p=p$ gives probabilism. This includes strictly proper scores that are not additive, as in Pettigrew~\cite{Pettigrew2022}, and permits infinite losses. It also applies to forecasts for a finite collection $\mathcal F\subseteq2^\Omega$ of events, for arbitrary $\Omega$: replace worlds by the distinct vectors $v_w=(\mathbf1_A(w))_{A\in\mathcal F}$ and put $c_p=\sum_wp_wv_w$. These are exactly the coherent forecasts. Continuous, strictly proper component scores give the unique Bayes forecasts, recovering the ordinary-dominance consequence of Predd et al.~\cite{Predd2009}. Those results also establish stronger forms of dominance. For merely proper scores, uniqueness can fail: incoherent forecasts can share coherent loss profiles and remain undominated~\cite{SSK2009}. {\color{red}Then no countable family of assignments satisfies both {\color{red}(S)} and {\color{red}(T)}, since by Theorem~\ref{thm:argument}(a) they would make every incoherent forecast dominated. When the score is continuous, $c_p=p$ still gives {\color{red}(S)}, as in the proof of Lemma~\ref{lem:loss-selection}; {\color{red}the failure lies in (T).}}

\subsection{Probabilism on arbitrary domains}
Let $\M_L=[0,1]^{2^\Omega}$ and let $\C$ be the functions $q:2^\Omega\to[0,1]$ such that $q(\Omega)=1$ and $q(A\cup B)=q(A)+q(B)$ whenever $A\cap B=\varnothing$. For $w\in\Omega$, the probability $\delta_w$ is defined by $\delta_w(A)=1$ if $w\in A$, and $0$ otherwise. These probabilities are given the coarsest topology making $q\mapsto q(A)$ continuous for every $A\subseteq\Omega$. For a simple function $h=\sum_i a_i\mathbf1_{A_i}$ on a finite partition $(A_i)$ of $\Omega$, put $E_qh=\sum_i a_iq(A_i)$. For a bounded function $h$, define $E_qh$ by uniformly approximating $h$ by simple functions; the limit is independent of the approximation. For $f\ge0$, define $E_qf=\sup\{E_qh:0\le h\le f,\ h\text{ bounded}\}$. Call $L$ \emph{quasi-strictly proper} if
\[
E_qL(q,\cdot)\le E_qL(m,\cdot)\quad(q\in\C,\ m\in\M_L),
\]
strictly whenever $m\notin\C$.
\begin{proposition}\label{prop:short-infinite}
If $L$ is quasi-strictly proper and each $q\mapsto L(q,\omega)$ is lower semicontinuous, the undominated models are exactly the finitely additive probabilities, and every other model is dominated by one.
\end{proposition}
\begin{proof}
First suppose $m$ dominates $q\in\C$, strictly improving at $w$. Quasi-strict propriety gives $E_qL(m)=E_qL(q)<\infty$; finiteness follows by comparison with any incoherent credence. Put $q_n=(1-1/n)q+\delta_w/n$. Propriety at $q_n$ and $q$ gives
\[
\tfrac1nL(q_n,w)+(1-\tfrac1n)E_qL(q)
\le E_{q_n}L(q_n)\le E_{q_n}L(m)
=\tfrac1nL(m,w)+(1-\tfrac1n)E_qL(q).
\]
Hence $L(q_n,w)\le L(m,w)$, contradicting lower semicontinuity at $q$. So $\C_L=\C$. Quasi-strict propriety excludes an incoherent model tied with a probability, since their $q$-expected losses would be equal. {\color{red}Thus $\C_L$ is closed under ties.

Apply Lemma~\ref{lem:loss-selection} with $\Theta\coloneqq\C$, compact as a closed subset of $[0,1]^{2^\Omega}$, and $c_q\coloneqq q$. For every finite nonempty $F$, let $\theta_F(p)$ be the probability with masses $p_\omega$ at the worlds $\omega\in F$; this map is continuous, since its value on an event $A$ is $\sum_{\omega\in F\cap A}p_\omega$. Quasi-strict propriety makes $E_pL(p)$ finite, so the Bayes inequality gives, for every model $m$, some $\omega\in\supp p$ with $L(\theta_F(p),\omega)\le L(m,\omega)$.}
\end{proof}
Nielsen~\cite{Nielsen2023} assumes continuity of $q\mapsto E_pL(q)$ for every probability $p$, together with local boundedness of coherent loss profiles, and obtains improvement at every world. The proposition requires only worldwise lower semicontinuity and drops boundedness, but concludes ordinary dominance. It does not claim Nielsen's stronger conclusion under weaker hypotheses.

\subsection{Conditional updating}
Let $\Omega$ be finite and let $E_1,\ldots,E_k$ be nonempty pairwise disjoint subsets with union $\Omega$. The models are strategies $(c,a_1,\ldots,a_k)\in([0,1]^{2^\Omega})^{k+1}$, specifying current credences and credences to adopt on learning each cell. Call it coherent when all these credences are probabilities,
\[
a_j(E_j)=1,\qquad c(A\cap E_j)=c(E_j)a_j(A)\quad\text{for every }A,j.
\]
Thus even a plan for a probability-zero cell must be probabilistic and concentrate on that cell. Let $J$ be a finite loss for credence functions on $2^\Omega$, strictly proper in the sense that each probability uniquely minimises its own expected loss, and lower semicontinuous on probabilities at every world, and assess strategies by
\[
L((c,a),\omega)=\alpha J(c,\omega)+\beta J(a_j,\omega)
\quad(\omega\in E_j),\qquad \alpha,\beta>0.
\]
\begin{proposition}\label{prop:short-updating}
The undominated strategies are exactly the coherent ones, and every incoherent strategy is dominated by a coherent one.
\end{proposition}
\begin{proof}
A weak dominator of a coherent strategy first minimises its $c$-expected loss, forcing its current credences to equal $c$ by strict propriety. Cancelling that common loss and taking expectations under each $a_j$ forces every conditional report to equal $a_j$. {\color{red}Thus $\C_L=\C$, and $\C_L$ is closed under ties, since a strategy tied with a coherent one weakly dominates it.}

For selection, put $\eta=1/(2k)$ and fix $u_j\in\Delta$ with $u_j(E_j)=1$ for each $j$. Given $p\in\Delta$, set
\[
r_j=\frac{p|_{E_j}+(\eta-p(E_j))_+u_j}{\max\{p(E_j),\eta\}},\qquad
\lambda_j=\frac{(p(E_j)-\eta)_+}{\sum_i(p(E_i)-\eta)_+},\qquad
q=\sum_j\lambda_jr_j,
\]
where $x_+=\max\{x,0\}$ and $p|_{E_j}(\omega)=p_\omega$ for $\omega\in E_j$, and $0$ otherwise. The denominator defining $\lambda_j$ is positive, since otherwise $1=\sum_jp(E_j)\le k\eta=1/2$. These define a continuous coherent strategy $(q,r_1,\ldots,r_k)$, Bayes for $q$, with $\supp q\subseteq\supp p$. Its Bayes inequality therefore gives, for every model $m$, a world $\omega\in\supp p$ with $L((q,r_1,\ldots,r_k),\omega)\le L(m,\omega)$. Its losses are lower semicontinuous in $p$, being sums of compositions of the lower semicontinuous $J$ with continuous probability maps. {\color{red}Apply Lemma~\ref{lem:loss-selection} with $\Theta\coloneqq\Delta$, $c_p\coloneqq(q,r_1,\ldots,r_k)$, and, for each nonempty $F\subseteq\Omega$, $\theta_F(p)_\omega=p_\omega$ for $\omega\in F$ and $\theta_F(p)_\omega=0$ otherwise.}
\end{proof}
The selected strategy need not be Bayes for the indexing prior $p$: a Bayes probability whose support is contained in $\supp p$ suffices. This avoids discontinuity of conditional probabilities at null cells. The result extends the finite-valued setting of Nielsen~\cite{Nielsen2021} without assuming that the score is additive or continuous on all forecasts. It also isolates the null-cell issue in Briggs and Pettigrew~\cite{Briggs2020}: merely conditionalizing on positive-probability cells does not ensure undominatedness. Invalid plans on null cells can be dominated without improvement at every world. Finiteness of $J$ matters: logarithmic losses can make such plans tied with coherent strategies. {\color{red}When that coherent strategy is undominated, so is the plan, and no countable family of assignments satisfies both {\color{red}(S)} and {\color{red}(T)}, by Theorem~\ref{thm:argument}(a).}

\subsection{Why coherent admissibility is insufficient}
Joyce's~\cite{Joyce2009} truth-directedness requires loss to decrease whenever every coordinate of a credence moves towards its truth value, with at least one strict improvement. Together with continuity, symmetry and coherent admissibility, this still does not ensure dominance of incoherent credences. For two worlds, take $\M_L=[0,1]^2$, $\C=\Delta$ and
\[
L_1(x)=(1-x_1+x_2)^2,\qquad L_2(x)=(1+x_1-x_2)^2.
\]
These losses have all four properties. Writing $d=x_1-x_2$, every $x$ is tied with the probability $((1+d)/2,(1-d)/2)$; distinct profiles are incomparable. Thus no model is dominated. Here $E_pL(x)=p_1(1-d)^2+p_2(1+d)^2$ is minimised at $d=p_1-p_2$. {\color{red}Thus $c_p=p$ is Bayes for $p$, and, as in the proof of Lemma~\ref{lem:loss-selection}, continuous losses verify {\color{red}(S)} with a single assignment. Evaluability and coherent admissibility hold, since no model is dominated. But every incoherent $x$ is tied with a probability and undominated, so {\color{red}(T)} fails for every countable family of assignments, by Theorem~\ref{thm:argument}(a).} Coherent admissibility alone cannot exclude incoherent models that perform exactly as well as coherent ones.


\subsection{Chance dominance}\label{sec:chance}
Pettigrew~\cite[Chapter~10]{Pettigrew2016} compares expected rather than actual accuracy. Let $W$ be finite, $\mathcal K\subseteq\Delta_W$ nonempty, and $\Theta=\overline{\operatorname{conv}\mathcal K}$, the Euclidean closure of the finite convex combinations of elements of $\mathcal K$. The members of $\mathcal K$ represent possible chances conditional on the agent's evidence. Let $\M_L=[0,1]^{\mathcal F}$ for a set of events $\mathcal F\subseteq2^W$, identify each probability $q$ with the vector $(q(A))_{A\in\mathcal F}$, and let $\C$ be the credences induced by $\Theta$. Suppose $I:\M_L\times W\to[0,\infty)$ is finite, each probability uniquely minimises its own expected loss, and $\theta\mapsto I(\theta,w)$ is lower semicontinuous. Index comparisons by $q\in\mathcal K$, putting $c\preceq_q d$ exactly when $E_qI(c)\le E_qI(d)$.
\begin{proposition}\label{prop:short-chance}
The undominated credences are exactly $\C$, and every other credence is dominated by a member of $\C$.
\end{proposition}
\begin{proof}
If $m$ weakly dominates $\theta\in\Theta$, the affine continuous function $q\mapsto E_qI(m)-E_qI(\theta)$ is nonpositive on $\mathcal K$, hence on $\Theta$. At $q=\theta$, unique Bayes optimality forces $m=\theta$ as credence functions. {\color{red}Thus $\C_L=\C$, and $\C_L$ is closed under ties, since a tied model weakly dominates.

Apply Lemma~\ref{lem:loss-selection}, with the chance functions $q\in\mathcal K$ as worlds and loss $E_qI$: take the compact parameter space $\Theta$, $c_\theta\coloneqq\theta$. For each $q\in\mathcal K$, the loss $\theta\mapsto E_qI(c_\theta)$ is lower semicontinuous, since it is a finite nonnegative combination of the functions $\theta\mapsto I(c_\theta,w)$. For each finite nonempty $F\subseteq\mathcal K$, use the continuous map $\theta_F(p)\coloneqq\sum_{q\in F}p_qq$. Since
\[
\sum_{q\in F}p_qE_qI(c)=E_{\theta_F(p)}I(c),
\]
Bayes optimality of $c_{\theta_F(p)}$ supplies, for every $m$, some $q\in\supp p$ with $E_qI(c_{\theta_F(p)})\le E_qI(m)$.}
\end{proof}
The chance set $\mathcal K$ may be infinite; additivity of the loss is not required. Pettigrew's Theorem~10.1.1 gives the stronger conclusion of improvement under every possible chance function, with an additional condition on the dominator's assessment of itself. Our conclusion is ordinary chance dominance by an undominated credence. Either argument requires a normative principle concerning chance dominance, not actual-loss dominance.


\bibliographystyle{plain}
\bibliography{ldg-references}
\end{document}
